\documentclass{article}

\usepackage[utf8]{inputenc}
%\usepackage[hebrew,english]{babel}
\usepackage[main=hebrew,english]{babel}
\usepackage{amssymb,amsmath,amsthm}
\usepackage{tikz}




\usepackage{pgfplots}
\pgfplotsset{compat=1.18}
\usepackage[margin=1in]{geometry}




\usepackage{float}















\usepackage{chngcntr}
\counterwithin{figure}{section}
\addto\captionshebrew{\renewcommand{\figurename}{איור}}




















% Hebrew-style theorem-like environments with section-based numbering

\newcounter{theoremHcounter}[section]
\renewcommand{\thetheoremHcounter}{\thesection.\arabic{theoremHcounter}}

\newcounter{defnHcounter}[section]
\renewcommand{\thedefnHcounter}{\thesection.\arabic{defnHcounter}}

\newcounter{excounter}[section]
\renewcommand{\theexcounter}{\thesection.\arabic{excounter}}

\newcounter{exscounter}[section]
\renewcommand{\theexscounter}{\thesection.\arabic{exscounter}}

\newcounter{exercounter}[section]
\renewcommand{\theexercounter}{\thesection.\arabic{exercounter}}

\newcounter{questcounter}[section]
\renewcommand{\thequestcounter}{\thesection.\arabic{questcounter}}

% Environments

\newenvironment{theoremH}
  {
    \par\noindent\refstepcounter{theoremHcounter}%
    \textbf{משפט~\thetheoremHcounter.}%
    \begin{itshape}
  }
  {
    \end{itshape}\par
  }

\newenvironment{prop}
  {
    \par\noindent\refstepcounter{theoremHcounter}%
    \textbf{טענה~\thetheoremHcounter.}%
    \begin{itshape}
  }
  {
    \end{itshape}\par
  }

\newenvironment{defnH}
  {
    \par\noindent\refstepcounter{defnHcounter}%
    \textbf{הגדרה~\thedefnHcounter.}%
    \begin{itshape}
  }
  {
    \end{itshape}\par
  }

\newenvironment{ex}
  {
    \par\noindent\refstepcounter{excounter}%
    \textbf{דוגמא~\theexcounter.}%
    \begin{itshape}
  }
  {
    \end{itshape}\par
  }

\newenvironment{exs}
  {
    \par\noindent\refstepcounter{exscounter}%
    \textbf{דוגמאות~\theexscounter.}%
    \begin{itshape}
  }
  {
    \end{itshape}\par
  }

\newenvironment{exer}
  {
    \par\noindent\refstepcounter{exercounter}%
    \textbf{תרגיל~\theexercounter.}%
    \begin{itshape}
  }
  {
    \end{itshape}\par
  }

\newenvironment{quest}
  {
    \par\noindent\refstepcounter{questcounter}%
    \textbf{שאלה~\thequestcounter.}%
    \begin{itshape}
  }
  {
    \end{itshape}\par
  }


















\begin{document}
\selectlanguage{hebrew}
\section{סדרות}
% ✅ Now using a real environment

\subsection{הגדרות ודוגמאות}
יתכן מאוד וכבר שמעתם על סדרות חשבוניות,
לדוגמא,
$1,2,3,4,5,\ldots$, 
וסדרות הנדסיות, 
לדוגמא, 
$1,\frac{1}{2},\frac{1}{4},\frac{1}{8},\frac{1}{16},\ldots$. 
אלה הן שתי משפחות של סדרות של מספרים. 
\begin{defnH}\rm
{\it סדרה של מספרים ממשיים} 
היא פונקציה 
$f$ 
אשר תחום הגדרתה הוא קבוצת המספרים הטבעיים וטווחה הוא קבוצת המספרים הממשיים. 
נסכים לכתוב 
$a_n$ 
במקום 
$f(n)$. 
\end{defnH}
\[
\begin{array}{cccccc}
1 & 2 & 3 & \ldots & n & \ldots \\
\downarrow & \downarrow & \downarrow & \downarrow & \downarrow & \downarrow \\
f(1) & f(2) & f(3) & \ldots & f(n) & \ldots
\end{array}
\]
שימו לב כי לכל סדרה אינסוף איברים. 

כמו לכל פונקציה, גם לסדרה יש הצגה גרפית. לכל מספר סידורי 
$n$ 
על ציר ה- 
$x$ 
נתאים 
נקודה 
$(n,a_n)$. 
לדוגמא, הסדרה שמתוארת באיור היא 
$a_n=\frac{1}{n}$. 


\vspace{1cm}

% Plot
%\selectlanguage{english}
%\begin{figure}[h]
%\centering
%\begin{tikzpicture}
%  \begin{axis}[
%    axis lines = left,
%    xlabel = {$n$},
%    ylabel = {$a_n=\frac{1}{n}$},
%    xtick = {1,2,3,4,5,6,7},
%    ytick = {0, 0.2, 0.4, 0.6, 0.8, 1.0},
%    ymin=0, ymax=1.1,
%    xmin=0.5, xmax=7.5,
%    width=12cm,
%    height=6cm,
%    grid=both,
%  ]
%    \addplot[
%      only marks,
%      mark=*,
%      mark size=2pt,
%      blue  % ← this makes the points blue
%    ] coordinates {
%      (1,1)
%      (2,0.5)
%      (3,0.333)
%      (4,0.25)
%      (5,0.2)
%      (6,0.166)
%      (7,0.142)
%    };
%  \end{axis}
%\end{tikzpicture}
%\selectlanguage{hebrew}
%\caption{שבעת האיברים הראשונים של $a_n=\frac{1}{n}$ על המישור הממשי}
%\label{fig:sequence-on-real-plane}
%\end{figure}



































































\selectlanguage{english}

\begin{figure}[H]
\centering
\begin{tikzpicture}
  \begin{axis}[
    axis lines=middle,
    xlabel={}, % Disable default xlabel
    ylabel={$a_n$},
    xmin=0, xmax=45,
    ymin=-0.05, ymax=1.1,
    width=12cm,
    height=7cm,
    xtick={1,2,3,4,5,10,20,40,41,42,43,44},
    xticklabels={1,2,3,4,5,10,20,,,,,},
    ytick={1/20, 0.1, 0.5, 1},
    yticklabels={ ,0.1,0.5,1},
    ytick style={draw=none},
    grid=none,
    xmajorgrids=false,
    ymajorgrids=false,
  ]

    % Sequence points: a_n = 1/n
    \addplot[
      only marks,
      mark=*,
      mark size=2pt,
      color=blue
    ] coordinates {
      (1,1)
      (2,0.5)
      (3,0.333)
      (4,0.25)
      (5,0.2)
      (10,0.1)
      (20,0.05)
      (40,0.025)
      (41,0.0244)
      (42,0.0238)
      (43,0.0233)
      (44,0.0227)
    };

    % Dashed red line at y=0 (the limit)
    \addplot[domain=0:45, samples=2, red, dashed, thick] {0};

    % Dashed tick mark on y-axis at y=0
    \draw[dashed] (axis cs:0,0) -- ++(-5pt,0);

    % Move label "n" 2mm toward the origin
    \node[xshift=-2mm] at (axis cs:45,-0.04) {$n$};

  \end{axis}
\end{tikzpicture}
\selectlanguage{hebrew}
\caption{כמה איברים ראשונים של הסדרה $a_n=\frac{1}{n}$}
\label{fig:limit-sequence}
\end{figure}

\selectlanguage{hebrew}






























































































\selectlanguage{hebrew}
ניתן להציג 
"גרף" אחר של סדרה. במקום לשרטט נקודות 
$(n,a_n)$ 
במישור, 
נשרטט נקודות 
$a_n$ 
על הציר הממשי. 
לשיטה הזאת יש חסרון - אחרי ששרטטנו את הנקודות, בלתי אפשרי להבין מהו הסדר בין האיברים האלה. 

% Plot
\selectlanguage{english}
\begin{figure}[h]
\centering
\begin{tikzpicture}[scale=10]
  \draw[->] (0,0) -- (1.1,0) node[below right] {$\mathbb{R}$};

  \foreach \n in {1,...,7} {
    \pgfmathsetmacro{\an}{1/\n}
    \filldraw[blue] ({\an}, 0) circle (0.3pt);
    \node[below] at ({\an}, 0) {$\frac{1}{\n}$};
  }

  \node[below] at (0,0) {$0$};
\end{tikzpicture}
\selectlanguage{hebrew}
\caption{שבעת האיברים הראשונים של $a_n=\frac{1}{n}$ על הציר הממשי}
\label{fig:sequence-on-real-line}
\end{figure}

\selectlanguage{hebrew}
נהוג לסמן סדרה ע"י 
$(a_n)_{n=1}^\infty$, 
כש- 
$a_{17}$ 
הוא האיבר ה- 
$17$ 
שלה ו- 
$a_{129}$ 
הוא האיבר ה- 
$129$ 
שלה. 
לעתים בכתיבת איברי הסדרה מוותרים על הסוגריים ורושמים פשוט 
$a_1,a_2,a_3,\ldots$. 

\begin{ex}\rm
$(a_n=n)_{n=1}^\infty$ 
היא סדרה של מספרים טבעיים המופיעים לפי הסדר ה"טבעי" שלהם, כלומר,
$$1,2,3,4,5,\ldots.$$
\end{ex}
\begin{ex}\rm
$(a_n=2n-1)_{n=1}^\infty$ 
היא סדרה של מספרים טבעיים אי-זוגיים, 
כלומר, 
$$1,3,5,7,9,\ldots.$$
\end{ex}
\begin{ex}\rm
הסדרה 
$a_n=1$ 
לכל 
$n$, 
כלומר 
$$1,1,1,1,1,\ldots.$$
\end{ex}
\begin{ex}\rm
$(a_n=(-1)^n)_{n=1}^\infty$ 
היא סדרה של מספרים 
$-1$ 
ו- 
$1$ 
המתחלפים, כלומר, 
$$-1,1,-1,1,-1,\ldots.$$
\end{ex}
\begin{ex}\rm
$\left(a_n=\frac{(-1)^{n+1}}{n}\right)_{n=1}^\infty$ 
היא הסדרה הבאה 
$$1,-\frac{1}{2},\frac{1}{3},-\frac{1}{4},\frac{1}{5},\ldots.$$
\end{ex}


קיימים אופנים שונים להגדרת סדרה. בכל אחת מהדוגמאות לעיל, סדרה הוגדרה ע"י נוסחא לחישוב איברה הכללי. 
זהו לא האופן היחיד. 
אחת השיטות הנפוצות להגדרת סדרה היא ע"י נוסחת נסיגה -  
השיטה שמאפשרת חישוב של איברה הכללי של סדרה דרך האיברים הקודמים. בשיטה זאת נחוצים 
{\it תנאיי התחלה}. 

\begin{ex}\rm
    תהא 
    $(a_n)_{n=1}^\infty$ 
    המתקבלת ע"י 
    $a_{n+1}=2a_n$ 
    לכל 
    $n$, 
    כאשר
    $a_1=1$. 
    נחשב כמה מאיבריה הראשונים של 
    $(a_n)_{n=1}^\infty$. 
    כאמור, 
    $a_1=1$. 
    מתקיים 
    $$a_2=2a_1=2\cdot 1=2,\ a_3=2a_2=2\cdot 2=4,\ a_4=2a_3=2\cdot4=8,\ldots.$$
    האם ניחשתם את הנוסחא לחישןב האיבר הכללי? 
    הכוכיחו )באינדוקציה מתמטית( כי צדקתם. 
\end{ex}

\begin{ex}\rm
תהא 
$(a_n)_{n=1}^\infty$ 
מוגדרת ע"י 
$a_{n}=a_{n-2}+a_{n-1}$ 
כאשר 
$a_1=1$
ו- 
$a_2=1$. 
)שימו לב כי סדרות עם נוסחאות נסיגה זהות שנבדלות בתנאיי התחילה הן סדרות שונות.(
לא קשה להשתכנע כי איבריה הראשונים של הסדרה הם
$1,1,2,3,5,8,13,21,\ldots$. 
קיימת גם נוסחא לחישוב של איברה הכללי של הסדרה הזאת
$$a_n=\frac{(1+\sqrt{5})^n-(1-\sqrt{5})^n}{2^n\sqrt{5}}.$$
\end{ex}

זאת סדרה "ידוענית" בתחום המתמטיקה, היא נקראת 
{\it סדרת פיבונצ'י}, 
על כינויו של מתמטיקאי איטלקי מימי הביניים- לאונרדו פיזנו. 
%למרבה הנס, הנוסחא המסובכת הזאת אכן "מייצרת" מספרים שלמים, למרות שזה לא נראה כך. 
פיזנו, בספרו 
{\it ספר החשבוניה} 
מביא דוגמא הבאה. 
זוג ארנבים, זכר ונקבה, נולד בתחילת חודש ראשון. 
לכל זוג נדרש חודש אחד כדי להתבגר ולהתרבות. 
כל זוג מוגר יולד זוג חדש, זכר ונקבה, כל חודש. 
בהנחה שארנבים לא מתי, כמה זוגות יהיו לאחר 
$n$
חודשים? 
השתכנעו כי 
$a_n$ 
זה אכן מספר הזוגות אחרי 
$n$ 
חודשים, 
כאשר 
$n=3,4,5,6$. 





































\begin{exer}\rm 
רשמו חמישה איברים ראשונים של הסדרה 
    $(a_n)_{n=1}^\infty$ 
    כאשר 
    \begin{enumerate}
        \item $a_n=\frac{n^3-1}{n^2+1}$.
        \item $a_n=\frac{n^2}{2^n}$. 
        \item $a_n=1-\frac{(-1)^n}{n}$. 
    \end{enumerate}
    \end{exer}
\begin{exer}\rm מצאו נוסחאות לאיבר הכללי של הסדרות הבאות:
    \begin{enumerate}
        \item $1,0,1,0,1,0\ldots$. 
        \item $\frac{5}{3},-\frac{10}{9},\frac{20}{27},-\frac{40}{81},\ldots$.
        \item $\frac{1}{1\cdot 3},\frac{1}{3\cdot 5},\frac{1}{5\cdot 7},\ldots$.
    \end{enumerate}
    \end{exer}
\begin{exer}\rm 
    תארו בצורה גרפית 
    )בשני אופנים( 
    את הסדרות הבאות: 
    \begin{enumerate}
        \item $a_n=1+\frac{1}{n}$. 
        \item $a_n=\frac{1-3n}{n}$. 
        \item $\frac{1+(-1)^n}{n}$. 
    \end{enumerate}
    \end{exer}
\begin{exer}\rm  האם  
    $42$ 
    הוא איבר של הסדרה 
    $(n^2-5n)_{n=1}^\infty$? 
    אם כן, מהו מיקומו בסדרה?
ענו על אותה השאלה עם המספרים 
$36$, 
$-4$
ו- 
$-6$.
\end{exer}
\begin{exer}\rm 
הציגו בצורה רקורסיבית את הסדרות הבאות: 
\begin{enumerate}
    \item $a_n=n$. 
    \item $a_n=\frac{1}{n}$.
    \item $(a_n)_{n=1}^\infty$ 
    היא סדרה כך שכל איבר, החל מאיבר השלישי, הוא ממוצע חשבוני של שני איברים קודמים. 
    כמה תנאיי התחלה נחוצים כדי לקבוע סדרה ביחידות? 
    \end{enumerate}
    כמה תנאי התחלה הגדרתם? 
\end{exer}

נחזור לדיון בסדרות. 
אנו הולכים להגדיר משפחות חשובות של סדרות שבעתיד הקרוב תשחקנה תפקיד חשוב. 

\begin{defnH}\rm
סדרה 
$(a_n)_{n=1}^\infty$ 
תקרא 
{\it יורדת} 
אם 
$a_n\geq a_{n+1}$ 
לכל 
$n$. 

סדרה 
$(a_n)_{n=1}^\infty$ 
תקרא 
{\it עולה} 
אם 
$a_n\leq a_{n+1}$ 
לכל 
$n$. 

סדרה 
$(a_n)_{n=1}^\infty$ 
תקרא 
{\it מונוטונית} 
אם 
היא עולה או יורדת.
\end{defnH}
%במילים אחרות, סדרה יורדת זו סדרה שאיבר הבא לא גדול מקודמו. 

\begin{ex}\rm
הסדרה 
$a_n=\frac{1}{n}$ 
היא יורדת. 
אכן 
$a_n-a_{n+1}=\frac{1}{n}-\frac{1}{n+1}=\frac{1}{n(n+1)}\geq 0$. 
\end{ex}

בדוגמא הקודמת כדי להראות שמספר אחד קטן או שווה מאחר - הסתכלנו על סימן של הפרש בין שניהם. 
זוהי טכניקה נפוצה. 

\begin{quest}\rm
האם הסדרה 
$a_n=1$ 
היא יורדת? 
\end{quest}
%במילים אחרות, סדרה עולה זו סדרה שאיבר הבא לא קטן מקודמו. 
\begin{ex}\rm
הסדרה 
$a_n=\frac{3n-5}{n}$ 
היא עולה. 
אכן 
$a_{n+1}-a_{n}=\frac{3(n+1)-5}{n+1}-\frac{3n}{n}=\frac{3}{n(n+1)}\geq 0$. 
גם הסדרה הקבועה 
$a_n=1$ 
היא עולה. 
\end{ex}

\begin{defnH}\rm
סדרה 
$(a_n)_{n=1}^\infty$ 
תקרא 
{\it  חסומה מלעיל} 
)או מלמעלה(
אם קיים מספר ממשי 
$M$ 
כך ש- 
$a_n\leq M$ 
לכל~$n$. 

סדרה 
$(a_n)_{n=1}^\infty$ 
תקרא 
{\it  חסומה מלרע} 
)או מלמטה(
אם קיים מספר ממשי 
$m$ 
כך ש- 
$a_n\geq m$ 
לכל~$n$. 

סדרה 
$(a_n)_{n=1}^\infty$ 
תקרא 
{\it חסומה} 
אם היא חסומה גם מלעיל וגם מלרע. 
\end{defnH}

\begin{ex}\rm
הסדרה 
$a_n=\frac{1}{n}$ 
חסומה מלעיל ע"י מספר 
$1$
אכן 
$\frac{1}{n}\leq 1$. 
\end{ex}


\begin{ex}\rm
הסדרה 
$a_n=\frac{3n-5}{n}$ 
חסומה מלרע ע"י מספר 
$-2$
אכן 
$\frac{3n-5}{n}=3-\frac{5}{n}\geq 3-5=-2$. 
\end{ex}

\begin{exer}\rm
    תנו דוגמא  
    \begin{enumerate}
        \item לסדרה שעולה ולא חסומה מלעיל.
        \item לסדרה שעולה וחסומה מלעיל.
        \item לסדרה מונוטונית ולא חסומה.
        \item לסדרה גם עולה וגם יורדת. 
        \item לסדרה חסומה ולא מונוטונית.
        \item לסדרה לא עולה ולא יורדת.
    \end{enumerate}
\end{exer}








\section
{גבול של סדרה}

נתבונן בסדרה של 
קירובים עשרוניים של 
$\frac{1}{3}$, 
דהיינו 
$$0, 0.3, 0.33, 0.333, 0.3333
\ldots.$$
ככל שמיקמוהים של איברי הסדרה יגדלו, איבריה של הסדרה הנ"ל "יתקרבו" יותר ויותר ל- 
$\frac{1}{3}$. 
במקרה זה נומר שהסדרה מתכנסת וגבולה הוא 
$\frac{1}{3}$. 
ראו איור 
\ref{decimal approximation of 1/3}.

\selectlanguage{english}
\begin{figure}[h]
\centering
\begin{tikzpicture}
  \begin{axis}[
    axis lines=middle,
    xlabel={$n$},
    ylabel={$a_n$},
    xmin=0, xmax=45,
    ymin=0, ymax=0.6,
    width=12cm,
    height=7cm,
    xtick={1,2,3,4,5,10,20,40,41,42,43,44},
    xticklabels={1,2,3,4,5,10,20,,,,,},
    ytick={0,0.3,0.33,0.3333},
    yticklabels={},
    ytick style={draw=none},
    grid=none,
    xmajorgrids=false,
    ymajorgrids=false,
  ]
    % Points
    \addplot[
      only marks,
      mark=*,
      mark size=2pt,
      color=blue
    ] coordinates {
      (1,0)
      (2,0.3)
      (3,0.33)
      (4,0.36)
      (5,0.38)
      (10,0.41)
      (20,0.46)
      (40,0.47)
      (41,0.474)
      (42,0.479)
      (43,0.482)
      (44,0.484)
    };

    % Interrupted red dashed horizontal line at y=0.52 (break at x=20)
    \draw[red, thick, dashed] (axis cs:0,0.4935) -- (axis cs:18,0.4935);
    \draw[red, thick, dashed] (axis cs:22,0.4935) -- (axis cs:45,0.4935);

    % Label y=1/3 aligned with the dashed line, shifted slightly down
    \node[red, anchor=south, yshift=-9pt] at (axis cs:20,0.4935) {$y=\frac{1}{3}$};

    % Dashed tick mark on y-axis at y=1/3 with label slightly above and left
    \draw[dashed] (axis cs:0,0.3333) -- ++(-5pt,0);
    \node[anchor=south east, xshift=-5pt, yshift=3pt] at (axis cs:0,0.3333) {$\frac{1}{3}$};

  \end{axis}
\end{tikzpicture}
\selectlanguage{hebrew}
\caption{סדרת הקירובים העשרוניים של $\frac{1}{3}$}
\label{decimal approximation of 1/3}
\end{figure}


\
\selectlanguage{hebrew}
עתה נתבונן בסדרה 
$a_n=\frac{1}{n}$, 
ראו איור 
\ref{fig:limit-sequence}. 
ככל ש- 
$n$ 
יגדל, איבריה של הסדרה "יתקרבו" ל- 
$0$. 










מיד נעבור להגדרה פורמלית של המושג המרכזי - גבול. 
אך לפני זה נכיר אות יוונית 
$\varepsilon$ - 
אפסילון. 
במרבית ספרי חדו"א כשמוזכר 
$\varepsilon$ 
הכוונה היא למספר חיובי 
"קטן" 
לא ידוע. 


































































%עיקר העניין שלנו בספר זה יתמקד בסדרות שבמובן מסוים יש להן התנהגות "יציבה" - סדרות שעם התקדמות האינדקסים, איבריהן מצטופפים ליד איזשהו מספר. 
%המספר הזה יקרא 
%{\it גבול של סדרה}. 
%מושג הגבול הינו המושג המרכזי בחשבון דיפרנציאלי ואינטגרלי. כפי שנראה בהמשך, מושגים רבים בתחום מבוססים עליו. 

\begin{defnH}\label{limit of sequence}\rm
תהא 
$(a_n)_{n=1}^\infty$ 
סדרה. 
אומרים של- 
$(a_n)_{n=1}^\infty$ 
קיים גבול 
$L\in\mathbb{R}$ 
אם לכל 
$\varepsilon>0$ 
קיים 
$N\in\mathbb{N}$ 
כך ש- 
$|a_n-L|<\varepsilon$, 
לכל 
$n> N$. 
את העובדה הזאת נציין בכתיבה 
$$\displaystyle\lim_{n \to \infty}a_n=L$$ 
או 
$$a_n\xrightarrow[n\longrightarrow\infty]{}L.$$

אם 
$\displaystyle\lim_{n \to \infty}a_n=L$, 
אז נומר ש- 
$(a_n)_{n=1}^\infty$ 
{\it מתכנסת}
או 
{\it שואפת} 
ל- 
$L$. 
\end{defnH}

נתעכב על ההגדרה של גבול. נבין בהדרגה. 
\begin{itemize}
    \item 
{\it לכל $\varepsilon>0$}:
המשמעות היא שכל מה שהולכים לומר אחרי זה נכון לכל מספר חיובי ללא יוצא דופן. 
\item 
{\it קיים $N\in\mathbb{N}$}: 
יחד עם הביטוי הקודם זה אומר שלכל מספר חיובי 
$\varepsilon$
נצליח למצוא מספר טבעי 
$N$. 
\item 
{\it כך ש- $|a_n-L|<\varepsilon$ לכל $n>N$}: 
וזה אומר שאיברי הסדרה 
$(a_n)_{n=1}^\infty$ 
"קרובים" למספר
$L$ 
עד כדי 
$\varepsilon$ 
החל ממקום 
$N$, 
כלומר, המרחק בין כל איברי הסדרה שמיקומם הוא לכל הפחות 
$N+1$ 
לגבול 
$L$ 
הוא קטן מ- 
$\varepsilon$. 

גרפית זה אומר שכל האיברים 
$a_{N+1},a_{N+2},a_{N+3},\ldots$ 
צריכים להמצא ברצועה אופקית שתחומה בין הישרים 
$y=L-\varepsilon$ 
ו- 
$y=L+\varepsilon$. 
ראו איור 
\ref{fig:limit-alternating-full}. 
\end{itemize}
ראו איור
\ref{fig:limit-alternating-full} 
שוב. 
ל- 
$\varepsilon>0$ 
שבאיור, 
$N$ 
שיתאים הוא 
$N=57$. 
אכן, אם תסתכלו טוב, 
$|a_{57}-L|\geq\varepsilon$, 
אבל 
$|a_n-L|<\varepsilon$ 
ל- 
$n=58,59,60,\ldots$. 
באיור אנו רואים שגם 
$|a_1-L|<\varepsilon$ 
וגם 
$|a_2-L|<\varepsilon$, 
אך 
$N=1$ 
ו- 
$N=2$ 
לא יספקו אותנו, מאחר ו -
$|a_3-L|\geq\varepsilon$. 
















\selectlanguage{english}
\begin{figure}[H]
\centering
\begin{tikzpicture}
  \begin{axis}[
    axis lines=middle,
    xlabel={$n$},
    ylabel={$a_n$},
    xmin=0, xmax=45,
    ymin=0.8, ymax=1.2,
    width=12cm,
    height=7cm,
    xtick={1,2,3,4,5,10,11,20,40,41,42,43,44},
    xticklabels={1,2,3,4,5,,58,128,},
    ytick={0.9,1,1.1},
    yticklabels={$L-\varepsilon$, $L$, $L+\varepsilon$},
    ytick style={draw=none},
    grid=none,
    xmajorgrids=false,
    ymajorgrids=false,
  ]

    % Points from n=1 to n=44
    \addplot[
      only marks,
      mark=*,
      mark size=2pt,
      color=blue
    ] coordinates {
      (1, 1.05)
      (2, 0.95)
      (3, 1.135)
      (4, 0.85)
      (5, 1.17)
      (6, 1.1666)
      (7, 0.8571)
      (8, 1.125)
      (9, 0.8888)
      (10, 1.1)
      (11, 0.9091)
      (12, 1.0833)
      (13, 0.9231)
      (14, 1.0714)
      (15, 0.9333)
      (16, 1.0625)
      (17, 0.9411)
      (18, 1.0556)
      (19, 0.9474)
      (20, 1.05)
      (40, 1.025)
      (41, 0.9756)
      (42, 1.0238)
      (43, 0.9767)
      (44, 1.0227)
    };

    % Dashed red line at y=1 (limit)
    \addplot[domain=0:45, samples=2, red, dashed, thick] {1};
    \node[red, anchor=south] at (axis cs:22,1) {%$y=1$
    };
    \draw[dashed] (axis cs:0,1) -- ++(-5pt,0);

    % Dashed red lines at y = 1.1 and y = 0.9
    \addplot[domain=0:45, samples=2, red, dashed] {1.1};
    \addplot[domain=0:45, samples=2, red, dashed] {0.9};

    % Labels for epsilon lines
    \node[red, anchor=south] at (axis cs:22,1.1) {%$y=L+\varepsilon$
    };
    \node[red, anchor=north] at (axis cs:22,0.9) {%$y = L - \varepsilon$
    };

  \end{axis}
\end{tikzpicture}
\selectlanguage{hebrew}
\caption{גבול של סדרה}
\label{fig:limit-alternating-full}
\end{figure}
\selectlanguage{hebrew}


















































%\vspace*{0.3cm}
%{\bf הערה מיותרת?} 
%בהגדרה הנ"ל אפשר לצמצם את טווח ה- 
%$\varepsilon$ 
%לכל 
%$\varepsilon<\varepsilon_0$ 
%עבור איזשהו 
%$\varepsilon_0>0$ 
%קבוע. 
%אכן, נניח כי אני מסוגלים "להפעיל" את הגדרת הגבול לכל 
%$\varepsilon<\frac{1}{100}$, 
%כלומר,
%לכל 
%$\varepsilon<\frac{1}{100}$
%קיים 
%$N$ 
%כך ש- 
%$|a_n-L|<\varepsilon$ 
%לכל 
%$N>N$. 
%עתה נשאל את עצמנו, מה לדוגמא עושים עם 
%$\varepsilon=\frac{1}{3}$? 
%ננהוג באופן הבא: נמצא 
%$N\left(\frac{1}{100}\right)$, 
%אזי, לכל 
%$N>N\left(\frac{1}{100}\right)$, 
%מתקיים
%$|a_n-L|<\frac{1}{100}$, 
%על אחת כמה וכמה, 
%$|a_n-L|<\frac{1}{3}$ 
%לכל 
%$n>N\left(\frac{1}{100}\right)$,
%כלומר, 
%$N\left(\frac{1}{100}\right)$ 
%"משרת" את כל האפסילונים שגדולים מ- 
%$\frac{1}{100}$. 


\begin{ex}\rm 
נראה כי
    $\displaystyle\lim_{n \to \infty}\frac{5n-12}{2n+3}=\frac{5}{2}$. 
    נפנה להגדרה
    \ref{limit of sequence}. 
    נקבע 
    $\varepsilon>0$. 
    מטרתנו למצוא 
    $N\in\mathbb{N}$ 
    כך ש- 
    $\left|a_n-\frac{5}{2}\right|<\varepsilon$ 
    לכל 
    $n>N$. 
    כלומר, עלינו למצוא פתרון לאי-שיוויון 
    $$\left|\frac{5n-12}{2n+3}-\frac{5}{2}\right|<\varepsilon.$$ 
יש לומר שאנו לא מעוניינים בפתרון כלשהו אלא בפתרון בעל צורה מסוימת - כל הטעי החל מאיזשהו טבעי 
$N$ 
צריכים להוות פתרון לאי-שיוויון. 
כלומר, אנו מעוניינים במציאת 
    $N$ 
    כך שאם 
    $n>N$ 
    יתקיים 
    $\left|\frac{5n-12}{2n+3}-\frac{5}{2}\right|<\varepsilon$. 
מתקיים 
        $$\varepsilon>\left|\frac{5n-12}{2n+3}-\frac{5}{2}\right|=\left|\frac{39}{4n+6}\right|.$$ 
        מכיוון שהביטוי באגף ימין הוא חיובי, אי-שיוויון האחרון שקול ל- 
        $$\frac{39}{4n+6}<\varepsilon.$$
        שבתורו שקול ל- 
        $$\frac{4n+6}{39}>\frac{1}{\varepsilon},$$
        כלומר, 
        $$4n+6>\frac{39}{\varepsilon},$$
        מכאן 
$$n>\frac{1}{4}\left(\frac{39}{\varepsilon}-6\right).$$  
המספר 
$\frac{39}{\varepsilon}-6$
יכול להיות שלילי
)לדוגמא, אם 
$\varepsilon=7$(, 
לכן 
    ניקח 
    $N$ 
    להיות השלם הכי קטן שגדול מ- 
    $\max\left\{\frac{1}{4}\left(\frac{39}{\varepsilon}-6\right),0\right\}$. 
    נראה כי ה-
    $N$
    הזה משרת את מטרתנו. 
    אם 
    $\max\left\{\frac{1}{4}\left(\frac{39}{\varepsilon}-6\right),0\right\}=\frac{1}{4}\left(\frac{39}{\varepsilon}-6\right)$, 
    אז
    לכל 
    $n>\frac{1}{4}\left(\frac{39}{\varepsilon}-6\right)$ 
    מתקיים
    $$\left|\frac{5n-12}{2n+3}-\frac{5}{2}\right|=\left|\frac{39}{4n+6}\right|<\left|\frac{39}{4\frac{1}{4}\left(\frac{39}{\varepsilon}-6\right)+6}\right|=|\varepsilon|=\varepsilon.$$
    ואם 
    $\max\left\{\frac{1}{4}\left(\frac{39}{\varepsilon}-6\right),0\right\}=0$, 
    נסיק כי 
    $\frac{1}{4}\left(\frac{39}{\varepsilon}-6\right)\leq 0$, 
    כלומר 
    $\frac{39}{\varepsilon}\leq 6$, 
    ז"א 
    $\varepsilon\geq\frac{39}{6}$, 
    אז נסיק כי 
    $$\left|\frac{5n-12}{2n+3}-\frac{5}{2}\right|=\left|\frac{39}{4n+6}\right|<\left|\frac{39}{4\cdot0+6}\right|=\frac{39}{6}\leq\varepsilon.$$ 
    בשני המקרים הצלחנו למצוא 
    $N\in\mathbb{N}$ 
    כך ש- 
    $$\left|\frac{5n-12}{2n+3}-\frac{5}{2}\right|<\varepsilon$$
    לכל 
    $n>N$, 
    כלומר, קיימנו את הנדרש בהגדרת גבול. 
\end{ex}    
\begin{ex}\rm 
נחזור לדוגמא הקודמת ונראה שמספר אחר, לדוגמא 
$\frac{3}{2}$ - 
לא מהווה גבול של הסדרה. 
נקבע 
$\varepsilon>0$ 
וננסה למצוא 
$N\in\mathbb{N}$ 
כך ש- 
$$\left|\frac{5n-12}{2n+3}-\frac{3}{2}\right|<\varepsilon,$$
כלומר,
$$\left|\frac{4n-33}{4n+6}\right|<\varepsilon.$$
אי-שיוויון זה שקול ל- 
$$-\varepsilon<\frac{4n-33}{4n+6}<\varepsilon.$$
נשים לב כי 
$\frac{4n-33}{4n+6}-1-\frac{39}{4n+6}$, 
אזי 
$$-\varepsilon<\frac{4n-33}{4n+6}<\varepsilon$$
אמ"מ
$$-\varepsilon<1-\frac{39}{4n+6}<\varepsilon$$
אמ"מ
$$-1-\varepsilon<-\frac{39}{4n+6}<-1+\varepsilon$$
אמ"מ
$$\frac{1-\varepsilon}{39}<\frac{1}{4n+6}<\frac{1+\varepsilon}{39}$$
אמ"מ
$$\frac{39}{1+\varepsilon}<4n+6<\frac{39}{1-\varepsilon}$$
אמ"מ
$$\frac{39}{1+\varepsilon}-6<4n<\frac{39}{1-\varepsilon}-6$$
אמ"מ
$$\frac{39}{4+4\varepsilon}-\frac{3}{2}<n<\frac{39}{4-4\varepsilon}-\frac{3}{2}.$$
נשים לב כי אם קיימים פתרונות לאי-שיוויון שפתרנו הרי  שמספרם סופי. 
נזכר כי בהגדרת הגבול נדרש פתרון שמכיל את כל הטבעיים החל מאיזשהו טבעי מסוים, בפרט אינסוף פתרונות. 
אם נלך לכיוון של הפרכת קיום הגבול, נצטרך לספק 
$\varepsilon>0$ 
כך שלכל 
$N\in\mathbb{N}$ 
קיים 
$n<N$ 
אשר יקיים 
$$\left|\frac{5n-12}{2n+3}-\frac{3}{2}\right|\geq\varepsilon,$$
כלומר 
$$\left|\frac{4n-33}{4n+6}\right|\geq\varepsilon.$$
מסובך לעבוד עם ערכים מוחלטים, לכן ננסה "להיפטר" מערך מוחלט, לשם כך נשים לב כי המכנה הוא חיובי ואילו המונה הוא חיובי רק כאשר 
$n\geq 9$. 
לכן נניח כי 
$n\geq 9$ 
ונמשיך 
$$\left|\frac{5n-12}{2n+3}-\frac{3}{2}\right|=\frac{4n-33}{4n+6}=1-\frac{39}{4n+6}\geq 1-\frac{39}{4\cdot 9+6}=\frac{1}{14}.$$
איך התוצאה שהתקבלה ניתנת לפירוש? 
נבחר 
$\varepsilon<\frac{1}{14}$ 
ואז בשביל 
$0<\varepsilon$ 
זה לא נוכל למצוא 
$N$ 
כך ש- 
$$\left|\frac{5n-12}{2n+3}-\frac{3}{2}\right|<\varepsilon$$ 
לכל 
$n>N$, 
הרי שהראנו כי 
$$\left|\frac{5n-12}{2n+3}-\frac{3}{2}\right|\geq\frac{1}{14}>\varepsilon.$$
זה מוכיח כי 
$\frac{3}{2}$ 
לא מהווה גבול של הסדרה. 
\end{ex} 


\begin{ex}\rm 
בדוגמא זו נראה שאף מספר לא מהווה גבול של הסדרה 
$a_n=(-1)^n$. 
קודם נשלול ש- 
$L=-1$ 
הוא גבול של הסדרה. 
נניח בשלילה ש- 
$L=-1$ 
הוא גבול. זה אומר שלכל 
$\varepsilon>0$ 
קיים 
$N$ 
כך שכל איברי הסדרה 
$(a_n)_{n=1}^\infty$ 
שייכים לסביבת-
$\varepsilon$ 
של 
$L$ 
ל- 
$n=N+1,N+2,N+3,\ldots$. 
בפרט, ל-
$\varepsilon=\frac{1}{10}$ 
קיים 
$N$ 
כך ש- 
$$|a_n-(-1)|<\frac{1}{10}$$
לכל 
$n>N$. 
אבל לערכים הזוגיים של 
$n$ 
$$|a_n-(-1)|=|2|>\frac{1}{10},$$
קיבלנו סתירה. זה אומר ש- 
$L=-1$ 
לא גבול של הסדרה. 
באופן דומה מראים כי 
$L=-1$ 
לא מהווה גבול של הסדרה הנתונה. 
עתה נניח כי 
$L\neq\pm1$ 
הוא גבולה של 
$(a_n)_{n=1}^\infty$. 
נבחר 
$$\varepsilon<\min\{|1-L|,|-1-L|\}.$$
מתקיים 
$|a_n-L|>\varepsilon$, 
מאחר ו - 
$$|a_n-L|=|1-L|>\varepsilon$$
ל- 
$n$ 
זוגי ו- 
$$|a_n-L|=|-1-L|>\varepsilon$$ 
ל- 
$n$ 
אי-זוגי. 
נשיק של- 
$\varepsilon>0$ 
מסוים 
$|a_n-L|<\varepsilon$ 
לכל 
$n$, 
כלומר, עבור 
$\varepsilon$ 
מסוים זה לא קיים 
$N$ 
כך ש- 
$|a_n-L|<\varepsilon$ 
לכל 
$n>N$. 
נסיק כי לסדרה הנתונה לא קיים גבול. 
\end{ex} 

%\begin{ex}\rm 
%    $\lim_{n\longrightarrow\infty}\frac{5n^2-3}{3n^2+n}\neq1$. 
%    עלינו להראות כי קיים 
%    $\varepsilon>0$ 
%    כך שלכל 
%    $N\in\mathbb{N}$ 
%    קיים 
%    $n> N$ 
%    המקיים 
%    $\left|\frac{5n^2-3}{3n^2+n}-1\right|\geq\varepsilon$. 
%    נתחיל בחישובים 
%    $$\left|\frac{5n^2-3}{3n^2+n}-1\right|=\left|\frac{2n^2-n-3}{3n^2+n}\right|\underset{n\geq 2}{=}\frac{2n^2-n-3}{3n^2+n}.$$
%    נראה באינדוקציה כי 
%    $\frac{2n^2-n-3}{3n^2+n}>\frac{1}{2}$ 
%    לכל 
%    $n\geq 5$. 
%    עבור 
%    $n=5$ 
%    נקבל 
%    $$\frac{2\cdot5^2-5-3}{3\cdot5^2+5}=\frac{21}{40}>\frac{1}{2}.$$
%    נניח כי הטענה נכונה לאיזשהו 
%    $n\geq 5$, 
%    כלומר 
%    $\frac{2n^2-n-3}{3n^2+n}>\frac{1}{2}$. 
%    מתקיים 
%    \begin{align*}
%        \frac{2(n+1)^2-(n+1)-3}{3(n+1)^2+(n+1)}-\frac{2n^2-n-3}{3n^2+n}&=\ \frac{2n^2+3n-2}{3n^2+7n+4}-\frac{2n^2-n-3}{3n^2+n}\\
%        &=\ \frac{5n^2+23n+12}{(3n^2+n)(3n^2+7n+4)}\\
%        &>\ 0,
%    \end{align*}
%    כלומר 
%    $$ \frac{2(n+1)^2-(n+1)-3}{3(n+1)^2+(n+1)}>\frac{2n^2-n-3}{3n^2+n}>\frac{1}{2},$$
%    כנדרש. 
%    אזי 
%    $\varepsilon=\frac{1}{2}$ 
%    משרת את מטרתנו. 
%\end{ex}
\vspace{0.3 cm}
מה ראינו בשלוש דוגמאות האחרונות? 
בשתי דוגמאות הראשונות ראינו כי מספר מסוים - 
$\frac{5}{2}$ - 
הוא גבול של הסדרה הנתונה ואילו מספר אחר - 
$\frac{3}{2}$ - 
הוא לאו. 
בדוגמא שלישית ראינו סדרה שאין לה גבול. 
שלוש התוצאות האלה באופן טבעי מובילות למשפט הבא. 
\vspace{0.3 cm}
\begin{prop}\label{uniqueness of limit}\rm
אם קיים גבול 
$\displaystyle\lim_{n \to \infty}a_n$, 
אז הוא יחיד. 
\end{prop}

\begin{proof}
יהיו 
$L,M\in\mathbb{R}$ 
כך ש- 
$$\displaystyle\lim_{n \to \infty}a_n=L,\ \displaystyle\lim_{n \to \infty}a_n=M$$ 
ו- 
$L\neq M$. 
בלי הגבלת כלליות
)משתמשים בביטוי זה כאשר בהוכחה יש מספר מקרים והוכחה של כל אחד מהם דומה להוכחה של מקרה אחד(
נניח כי 
$L>M$
יהא 
$\varepsilon=\frac{L-M}{2}$. 
מהעובדה ש- 
$\displaystyle\lim_{n \to \infty}a_n=L$, 
קיים
$N_1\in\mathbb{N}$ 
כך ש- 
$|a_n-L|<\varepsilon$ 
לכל 
$n>N_1$, 
בפרט, 
$a_n>\frac{L+M}{2}$ 
לכל 
$n>N_1$. 
כמוכן, 
מהעובדה ש- 
$\displaystyle\lim_{n \to \infty}a_n=M$, 
קיים
$N_2\in\mathbb{N}$ 
כך ש- 
$|a_n-M|<\varepsilon$ 
לכל 
$n>N_2$, 
בפרט, 
$a_n>\frac{L+M}{2}$ 
לכל 
$n>N_2$. 
קיבלנו כי לכל 
$n> N=\max\{N_1,N_2\}$ 
מתקיים בו-זמנית 
$a_n>\frac{L+M}{2}$ 
ו- 
$a_n<\frac{L+M}{2}$, 
סתירה. 
\end{proof} 

\selectlanguage{english}
\begin{center}
\begin{tikzpicture}

  % Define L, M, and epsilon
  \def\L{6}
  \def\M{2}
  \def\eps{2} % (L - M)/2 = 2

  % Draw horizontal axis extended more to the right
  \draw[->] (-2,0) -- (10,0);

  % Mark L point with filled red dot and red label below
  \draw[fill=red] (\L,0) circle (2pt) node[below=6pt,red] {$\textcolor{red}{L}$};

  % Mark M point with filled blue dot and blue label above
  \draw[fill=blue] (\M,0) circle (2pt) node[above=6pt,blue] {$\textcolor{blue}{M}$};

  % Draw open circles at endpoints and label:
  % Left interval endpoints above
  \foreach \x/\label in {
    \L-\eps/{$L - \varepsilon$},
    \L+\eps/{$L + \varepsilon$}
  } {
    \draw (\x,0) circle (3pt);
    \node[above=6pt] at (\x,0) {\small \label};
  }
  % Right interval endpoints below
  \foreach \x/\label in {
    \M-\eps/{$M - \varepsilon$},
    \M+\eps/{$M + \varepsilon$}
  } {
    \draw (\x,0) circle (3pt);
    \node[below=6pt] at (\x,0) {\small \label};
  }

  % Draw point close to L inside left interval with orange color and label above
  \def\pointL{5.7}
  \draw[fill=orange] (\pointL,0) circle (2pt) node[above=6pt,orange] {$a_{N+1}$};

  % Draw point close to M inside right interval with orange color and label below
  \def\pointM{2.5}
  \draw[fill=orange] (\pointM,0) circle (2pt) node[below=6pt,orange] {$a_{N+1}$};

\end{tikzpicture}
\end{center}

\selectlanguage{hebrew}
בהוכחת הטענה 
\ref{uniqueness of limit}
ראינו ברמיזה עובדה חשובה שתנוסח בטענה הבאה. 

\begin{prop}\rm
אם 
$\lim_{n\longrightarrow \infty} a_n=L$, 
אז כל סביבה של 
$L$ 
מכילה 
כמעט כל איברי הסדרה
$(a_n)_{n=1}^\infty$ 
\newline
)ביתר הדיוק, מכילה את כל האיברים של הסדרה החל ממקום מסוים בפרט סדרה שמתכנסת
\newline
היא חסומה.(
\end{prop}
\begin{proof}
ע"פ ההגדה, ל- 
$\varepsilon>0$ 
הנתון מראש 
קיים 
$N\in\mathbb{N}$ 
כך ש-  
$|a_n-L|<\varepsilon$ 
לכל 
$n>N$, 
כלומר 
$$L-\varepsilon<a_n<L+\varepsilon,$$
זאת אומרת שמחוץ לסביבה
$(L-\varepsilon,L+\varepsilon)$ 
יכולים להימצא רק האיברים 
$a_1,a_2,\ldots,a_{N-1}$.
ניקח עתה 
$$m=\min\{L-\varepsilon,a_1,a_2,\ldots,a_{N}\},\ M=\max\{L+\varepsilon,a_1,a_2,\ldots,a_{N}\},$$ 
אזי 
$m\leq a_n\leq M$ 
לכל 
$n\in\mathbb{N}$, 
כלומר 
$(a_n)_{n=1}^\infty$ 
היא סדרה חסומה. 
\end{proof}

\vspace*{0.3cm}


\subsection{ משפטי התכנסות}

בהינתן שתי סדרות
$(a_n)_{n=1}^\infty$
ו- 
$(b_n)_{n=1}^\infty$
נגדיר 
\begin{itemize}
    \item 
{\it סכום} 
שלהם להיות סדרה שאיברה ה- 
$k$ 
הוא סכום של איבר ה- 
$k$ 
של 
$(a_n)_{n=1}^\infty$
עם איבר ה- 
$k$ 
של 
$(b_n)_{n=1}^\infty$, 
דהיינו, 
$(a_n+b_n)_{n=1}^\infty$.  
    \item 
{\it הפרש} 
בין 
$(a_n)_{n=1}^\infty$
ל- 
$(b_n)_{n=1}^\infty$ 
להיות סדרה שאיברה ה- 
$k$ 
הוא הפרש בין איבר ה- 
$k$ 
של 
$(a_n)_{n=1}^\infty$
ואיבר ה- 
$k$ 
של 
$(b_n)_{n=1}^\infty$, 
דהיינו, 
$(a_n-b_n)_{n=1}^\infty$.
    \item 
{\it מכפלה} 
של  
$(a_n)_{n=1}^\infty$
ב- 
$(b_n)_{n=1}^\infty$ 
להיות סדרה שאיברה ה- 
$k$ 
הוא מכפלת איבר ה- 
$k$ 
של 
$(a_n)_{n=1}^\infty$
באיבר ה- 
$k$ 
של 
$(b_n)_{n=1}^\infty$, 
דהיינו, 
$(a_nb_n)_{n=1}^\infty$.
\item
{\it מנה} 
בין 
$(a_n)_{n=1}^\infty$
ל-  
$(b_n)_{n=1}^\infty$ 
להיות סדרה שאיברה ה- 
$k$ 
הוא חלוקת איבר ה- 
$k$ 
של 
$(a_n)_{n=1}^\infty$
באיבר ה- 
$k$ 
של 
$(b_n)_{n=1}^\infty$, 
דהיינו, 
$\left(\frac{a_n}{b_n}\right)_{n=1}^\infty$. 
להבדיל מההגדרות הקודמות, בהגדרה זו קיימת דרישה נוספת - 
$b_n\neq 0$ 
לכל 
$n$. 
\end{itemize}

המשפט הבא, הידוע בשם 
{\it אריתמטיקת הגבולות},
 מגלה קשר בין גבולות של סכום, הפרש, מכפלה ומנה אם לגבולות של 
 $(a_n)_{n=1}^\infty$ 
 ו- 
 $(b_n)_{n=1}^\infty$. 





\vspace*{0.3cm}
\begin{theoremH}\rm
יהיו  
$(a_n)_{n=1}^\infty$
ו- 
$(b_n)_{n=1}^\infty$ 
סדרות
נניח כי קיימים הגבולות 
$$\displaystyle\lim_{n \to \infty}a_n=L_1,\ 
\displaystyle\lim_{n \to \infty}b_n=L_2.$$ 
אזי 
\begin{enumerate}
    \item[1.] הגבול 
    $\displaystyle\lim_{n \to \infty}(a_n\pm b_n)$ 
    קיים 
    ו- 
    $\displaystyle\lim_{n \to \infty}(a_n\pm b_n)=L_1\pm L_2$. 
    \item[2.] הגבול 
    $\displaystyle\lim_{n \to \infty}(a_n\cdot b_n)$ 
    קיים 
    ו- 
    $\displaystyle\lim_{n \to \infty}(a_n\cdot b_n)=L_1\cdot L_2$. 
    \item[3.] אם 
    $L_2\neq 0$, 
    אז קיים הגבול 
    $\displaystyle\lim_{n \to \infty}\frac{a_n}{b_n}$ 
    ו- 
    $\displaystyle\lim_{n \to \infty}\frac{a_n}{b_n}=\frac{L_1}{L_2}$. 
\end{enumerate}
\end{theoremH}

התוצאות לא מפתיעות. נספק הסבר אינטואיטיבי לנכונות הטענה הראשונה במשפט. 
אם 
$\displaystyle\lim_{n \to \infty}(a_n)=L_1$ 
ו- 
$\displaystyle\lim_{n \to \infty}(b_n)=L_2$, 
אז, 
"עם הזמן",
איבריה של 
$(a_n)_{n=1}^\infty$ 
מתקרבים ל- 
$L_1$ 
ואיבריה של 
$(b_n)_{n=1}^\infty$ 
מתקרבים ל- 
$L_2$, 
לכן איבריה של 
$(a_n+b_n)_{n=1}^\infty$ 
מתקרבים ל- 
$L_1+L_2$. 
שאר הטענות מוסברות באופן דומה. עתה נעבור להוכחת המשפט. 

\end{document}